% Generated draft: requires mathematical and source review.
\documentclass[11pt]{article}
\usepackage[margin=25mm]{geometry}
\usepackage{fontspec}
\setmainfont{Latin Modern Roman}
\usepackage{amsmath,amssymb,mathtools,unicode-math}
\setmathfont{Latin Modern Math}
\usepackage{enumitem,needspace,xurl,hyperref}
\setlength{\parindent}{0pt}
\setlength{\parskip}{5pt}
\setlength{\emergencystretch}{3em}
\newcommand{\N}{\mathbb{N}}
\newcommand{\Z}{\mathbb{Z}}
\newcommand{\Q}{\mathbb{Q}}
\newcommand{\R}{\mathbb{R}}
\newcommand{\C}{\mathbb{C}}
\newcommand{\F}{\mathbb{F}}
\newcommand{\A}{\mathbb{A}}
\newcommand{\PP}{\mathbb{P}}
\newcommand{\E}{\mathbb{E}}
\newcommand{\eps}{\varepsilon}
\newcommand{\dd}{\,\mathrm{d}}
\begin{document}
\section*{UnsolvedMath ID: 165}
\textbf{Source problem code:} catalogue entry
\subsection{Definitions and mathematical statement}
For a positive integer $N$, write $[N]=\{1,\ldots,N\}$. Let $A\subset[N]$, let $c>0$ be a real parameter, and let $r\geq1$ be an integer parameter. Count a sum $x+y$ and the reversed sum $y+x$ as the same representation. With the usual convention that diagonal sums are also counted, this means
\[
r_A(n)=|\{(x,y)\in A\times A:x\leq y,\ x+y=n\}|\qquad(n\in\Z).
\]
If the intended convention excludes $x=y$, one must replace $x\leq y$ by $x<y$; the supplied qualification explicitly resolves ordering but not the treatment of diagonal sums.

The hypotheses in the source are
\[
|A|\geq c\sqrt N,\qquad \forall n\in\Z,\quad r_A(n)\leq r.
\]
The requested research task is to determine what structural information can be deduced about $A$ from these hypotheses. The source does not specify a structural conclusion or equivalence relation, how $c$ and $r$ are quantified, whether they remain fixed as $N\to\infty$, or the permitted dependence of any conclusion on these parameters. Thus the displayed assumptions are precise under the stated representation convention, but they do not yet define a theorem to prove.

Under the diagonal-counting convention, the case $r=1$ is the Sidon condition that each sum has at most one unordered representation. The supplied literature triage dated 2026-08-17 reports that Eberhard--Manners give a projective-plane mechanism underlying all known dense Sidon constructions, classify the Desarguesian cases within that framework, exhibit many non-Desarguesian examples, and supply smaller algebraic examples. It reports their claim that all dense Sidon sets arise through this mechanism as a conjecture, not a classification theorem for arbitrary $A$ satisfying the displayed hypotheses. The source's exact meaning of ``dense'' and the construction mechanism are not defined in this record.

The registry marks the problem open, while its dated triage uses partially solved to describe construction-level progress. The same triage says Green presents the task as a vague structural programme and regards even formulating a conjecture as problematic. In particular, the background's assertion that all such sets have special structure is not a proved universal conclusion.
\subsection{Short English statement}
What structure is forced when a subset of $\{1,\ldots,N\}$ has at least $c\sqrt N$ elements but every integer has at most $r$ unordered representations as a sum of two of its elements? The source leaves the structural target and parameter quantifiers unspecified; known projective-plane constructions do not classify every such set.
\subsection{Sources}
\begin{itemize}
\item Sean Eberhard and Freddie Manners, ``The apparent structure of dense Sidon sets,'' Electronic Journal of Combinatorics 30 (2023), no. 1, P1.33; arXiv:2107.05744: \url{https://arxiv.org/abs/2107.05744}.
\item Ben Green, 100 Open Problems, Problem 34, most recent update December 2025 according to the supplied record: \url{https://people.maths.ox.ac.uk/greenbj/papers/open-problems.pdf}. The supplied triage attributes the programme's explicitly vague framing to this reference. The citation is retained without opening or generating a PDF.
\item Local registry data/unsolvedmath/problems.json, permanent catalogue entry, exact source code catalogue entry, ``Structure of Sets with Bounded Representation,'' in Ben Green's 100 Open Problems; statement, background, and literature triage dated 2026-08-17. Dataset attribution: Ulam AI and the UnsolvedMath Contributors.
\end{itemize}
\end{document}
